% Options for packages loaded elsewhere
% Options for packages loaded elsewhere
\PassOptionsToPackage{unicode}{hyperref}
\PassOptionsToPackage{hyphens}{url}
\PassOptionsToPackage{dvipsnames,svgnames,x11names}{xcolor}
%
\documentclass[
  letterpaper,
  DIV=11,
  numbers=noendperiod]{scrartcl}
\usepackage{xcolor}
\usepackage{amsmath,amssymb}
\setcounter{secnumdepth}{-\maxdimen} % remove section numbering
\usepackage{iftex}
\ifPDFTeX
  \usepackage[T1]{fontenc}
  \usepackage[utf8]{inputenc}
  \usepackage{textcomp} % provide euro and other symbols
\else % if luatex or xetex
  \usepackage{unicode-math} % this also loads fontspec
  \defaultfontfeatures{Scale=MatchLowercase}
  \defaultfontfeatures[\rmfamily]{Ligatures=TeX,Scale=1}
\fi
\usepackage{lmodern}
\ifPDFTeX\else
  % xetex/luatex font selection
\fi
% Use upquote if available, for straight quotes in verbatim environments
\IfFileExists{upquote.sty}{\usepackage{upquote}}{}
\IfFileExists{microtype.sty}{% use microtype if available
  \usepackage[]{microtype}
  \UseMicrotypeSet[protrusion]{basicmath} % disable protrusion for tt fonts
}{}
\makeatletter
\@ifundefined{KOMAClassName}{% if non-KOMA class
  \IfFileExists{parskip.sty}{%
    \usepackage{parskip}
  }{% else
    \setlength{\parindent}{0pt}
    \setlength{\parskip}{6pt plus 2pt minus 1pt}}
}{% if KOMA class
  \KOMAoptions{parskip=half}}
\makeatother
% Make \paragraph and \subparagraph free-standing
\makeatletter
\ifx\paragraph\undefined\else
  \let\oldparagraph\paragraph
  \renewcommand{\paragraph}{
    \@ifstar
      \xxxParagraphStar
      \xxxParagraphNoStar
  }
  \newcommand{\xxxParagraphStar}[1]{\oldparagraph*{#1}\mbox{}}
  \newcommand{\xxxParagraphNoStar}[1]{\oldparagraph{#1}\mbox{}}
\fi
\ifx\subparagraph\undefined\else
  \let\oldsubparagraph\subparagraph
  \renewcommand{\subparagraph}{
    \@ifstar
      \xxxSubParagraphStar
      \xxxSubParagraphNoStar
  }
  \newcommand{\xxxSubParagraphStar}[1]{\oldsubparagraph*{#1}\mbox{}}
  \newcommand{\xxxSubParagraphNoStar}[1]{\oldsubparagraph{#1}\mbox{}}
\fi
\makeatother

\usepackage{color}
\usepackage{fancyvrb}
\newcommand{\VerbBar}{|}
\newcommand{\VERB}{\Verb[commandchars=\\\{\}]}
\DefineVerbatimEnvironment{Highlighting}{Verbatim}{commandchars=\\\{\}}
% Add ',fontsize=\small' for more characters per line
\usepackage{framed}
\definecolor{shadecolor}{RGB}{241,243,245}
\newenvironment{Shaded}{\begin{snugshade}}{\end{snugshade}}
\newcommand{\AlertTok}[1]{\textcolor[rgb]{0.68,0.00,0.00}{#1}}
\newcommand{\AnnotationTok}[1]{\textcolor[rgb]{0.37,0.37,0.37}{#1}}
\newcommand{\AttributeTok}[1]{\textcolor[rgb]{0.40,0.45,0.13}{#1}}
\newcommand{\BaseNTok}[1]{\textcolor[rgb]{0.68,0.00,0.00}{#1}}
\newcommand{\BuiltInTok}[1]{\textcolor[rgb]{0.00,0.23,0.31}{#1}}
\newcommand{\CharTok}[1]{\textcolor[rgb]{0.13,0.47,0.30}{#1}}
\newcommand{\CommentTok}[1]{\textcolor[rgb]{0.37,0.37,0.37}{#1}}
\newcommand{\CommentVarTok}[1]{\textcolor[rgb]{0.37,0.37,0.37}{\textit{#1}}}
\newcommand{\ConstantTok}[1]{\textcolor[rgb]{0.56,0.35,0.01}{#1}}
\newcommand{\ControlFlowTok}[1]{\textcolor[rgb]{0.00,0.23,0.31}{\textbf{#1}}}
\newcommand{\DataTypeTok}[1]{\textcolor[rgb]{0.68,0.00,0.00}{#1}}
\newcommand{\DecValTok}[1]{\textcolor[rgb]{0.68,0.00,0.00}{#1}}
\newcommand{\DocumentationTok}[1]{\textcolor[rgb]{0.37,0.37,0.37}{\textit{#1}}}
\newcommand{\ErrorTok}[1]{\textcolor[rgb]{0.68,0.00,0.00}{#1}}
\newcommand{\ExtensionTok}[1]{\textcolor[rgb]{0.00,0.23,0.31}{#1}}
\newcommand{\FloatTok}[1]{\textcolor[rgb]{0.68,0.00,0.00}{#1}}
\newcommand{\FunctionTok}[1]{\textcolor[rgb]{0.28,0.35,0.67}{#1}}
\newcommand{\ImportTok}[1]{\textcolor[rgb]{0.00,0.46,0.62}{#1}}
\newcommand{\InformationTok}[1]{\textcolor[rgb]{0.37,0.37,0.37}{#1}}
\newcommand{\KeywordTok}[1]{\textcolor[rgb]{0.00,0.23,0.31}{\textbf{#1}}}
\newcommand{\NormalTok}[1]{\textcolor[rgb]{0.00,0.23,0.31}{#1}}
\newcommand{\OperatorTok}[1]{\textcolor[rgb]{0.37,0.37,0.37}{#1}}
\newcommand{\OtherTok}[1]{\textcolor[rgb]{0.00,0.23,0.31}{#1}}
\newcommand{\PreprocessorTok}[1]{\textcolor[rgb]{0.68,0.00,0.00}{#1}}
\newcommand{\RegionMarkerTok}[1]{\textcolor[rgb]{0.00,0.23,0.31}{#1}}
\newcommand{\SpecialCharTok}[1]{\textcolor[rgb]{0.37,0.37,0.37}{#1}}
\newcommand{\SpecialStringTok}[1]{\textcolor[rgb]{0.13,0.47,0.30}{#1}}
\newcommand{\StringTok}[1]{\textcolor[rgb]{0.13,0.47,0.30}{#1}}
\newcommand{\VariableTok}[1]{\textcolor[rgb]{0.07,0.07,0.07}{#1}}
\newcommand{\VerbatimStringTok}[1]{\textcolor[rgb]{0.13,0.47,0.30}{#1}}
\newcommand{\WarningTok}[1]{\textcolor[rgb]{0.37,0.37,0.37}{\textit{#1}}}

\usepackage{longtable,booktabs,array}
\newcounter{none} % for unnumbered tables
\usepackage{calc} % for calculating minipage widths
% Correct order of tables after \paragraph or \subparagraph
\usepackage{etoolbox}
\makeatletter
\patchcmd\longtable{\par}{\if@noskipsec\mbox{}\fi\par}{}{}
\makeatother
% Allow footnotes in longtable head/foot
\IfFileExists{footnotehyper.sty}{\usepackage{footnotehyper}}{\usepackage{footnote}}
\makesavenoteenv{longtable}
\usepackage{graphicx}
\makeatletter
\newsavebox\pandoc@box
\newcommand*\pandocbounded[1]{% scales image to fit in text height/width
  \sbox\pandoc@box{#1}%
  \Gscale@div\@tempa{\textheight}{\dimexpr\ht\pandoc@box+\dp\pandoc@box\relax}%
  \Gscale@div\@tempb{\linewidth}{\wd\pandoc@box}%
  \ifdim\@tempb\p@<\@tempa\p@\let\@tempa\@tempb\fi% select the smaller of both
  \ifdim\@tempa\p@<\p@\scalebox{\@tempa}{\usebox\pandoc@box}%
  \else\usebox{\pandoc@box}%
  \fi%
}
% Set default figure placement to htbp
\def\fps@figure{htbp}
\makeatother





\setlength{\emergencystretch}{3em} % prevent overfull lines

\providecommand{\tightlist}{%
  \setlength{\itemsep}{0pt}\setlength{\parskip}{0pt}}



 


\KOMAoption{captions}{tableheading}
\makeatletter
\@ifpackageloaded{caption}{}{\usepackage{caption}}
\AtBeginDocument{%
\ifdefined\contentsname
  \renewcommand*\contentsname{Table of contents}
\else
  \newcommand\contentsname{Table of contents}
\fi
\ifdefined\listfigurename
  \renewcommand*\listfigurename{List of Figures}
\else
  \newcommand\listfigurename{List of Figures}
\fi
\ifdefined\listtablename
  \renewcommand*\listtablename{List of Tables}
\else
  \newcommand\listtablename{List of Tables}
\fi
\ifdefined\figurename
  \renewcommand*\figurename{Figure}
\else
  \newcommand\figurename{Figure}
\fi
\ifdefined\tablename
  \renewcommand*\tablename{Table}
\else
  \newcommand\tablename{Table}
\fi
}
\@ifpackageloaded{float}{}{\usepackage{float}}
\floatstyle{ruled}
\@ifundefined{c@chapter}{\newfloat{codelisting}{h}{lop}}{\newfloat{codelisting}{h}{lop}[chapter]}
\floatname{codelisting}{Listing}
\newcommand*\listoflistings{\listof{codelisting}{List of Listings}}
\makeatother
\makeatletter
\makeatother
\makeatletter
\@ifpackageloaded{caption}{}{\usepackage{caption}}
\@ifpackageloaded{subcaption}{}{\usepackage{subcaption}}
\makeatother
\usepackage{bookmark}
\IfFileExists{xurl.sty}{\usepackage{xurl}}{} % add URL line breaks if available
\urlstyle{same}
\hypersetup{
  pdftitle={Week-01: GitHub Workflow Tutorial},
  colorlinks=true,
  linkcolor={blue},
  filecolor={Maroon},
  citecolor={Blue},
  urlcolor={Blue},
  pdfcreator={LaTeX via pandoc}}


\title{Week-01: GitHub Workflow Tutorial}
\author{}
\date{}
\begin{document}
\maketitle

\renewcommand*\contentsname{Table of contents}
{
\setcounter{tocdepth}{3}
\tableofcontents
}

\subsection{\texorpdfstring{Working with \texttt{course-hub} and your
personal tutorial repository
\texttt{tutorials-exercises}.}{Working with course-hub and your personal tutorial repository tutorials-exercises.}}\label{working-with-course-hub-and-your-personal-tutorial-repository-tutorials-exercises.}

\textbf{Estimated time:} 60-75 minutes

\subsection{Background}\label{background}

In this course, you will use two different GitHub repositories for
in-class work (you will use others for homework assignments):

\begin{enumerate}
\def\labelenumi{\arabic{enumi}.}
\tightlist
\item
  \textbf{\texttt{course-hub}} contains course materials provided by the
  instructor.
\item
  \textbf{Your personal tutorial repository} contains the tutorials and
  exercises you complete during class.
\end{enumerate}

You will \textbf{pull} new materials from \texttt{course-hub}, but you
will complete, save, commit, and push your work in your personal
tutorial repository.

\begin{center}\rule{0.5\linewidth}{0.5pt}\end{center}

\subsection{Learning outcomes}\label{learning-outcomes}

By the end of this tutorial, you will be able to:

\begin{itemize}
\tightlist
\item
  distinguish between the central course repository and your personal
  tutorial repository;
\item
  explain the difference between a remote GitHub repository and a local
  repository on your computer;
\item
  clone a repository from GitHub to your computer;
\item
  pull updates from GitHub;
\item
  copy course materials into the correct repository;
\item
  review, stage, commit, and push changes;
\item
  confirm that your work appears on GitHub; and
\item
  reopen an existing RStudio Project without cloning it again.
\end{itemize}

\section{1. Create a course folder on your
computer}\label{create-a-course-folder-on-your-computer}

Call it \texttt{skill-building-in-R/}. This folder will house many
GitHub repositories - the course-hub, the tutorials-exercises (on your
personal account), and your homework repositories. For now, we will
focus on the course-hub and the tutorials-exercises repositories.

\texttt{skill-building-in-R/} is your parent folder.

\section{2. Understand the two
repositories}\label{understand-the-two-repositories}

\subsection{\texorpdfstring{Central course repository:
\texttt{course-hub}}{Central course repository: course-hub}}\label{central-course-repository-course-hub}

The instructor controls this repository. It contains template files for
the tutorials and exercises.

\begin{Shaded}
\begin{Highlighting}[]
\NormalTok{skill{-}building{-}in{-}R/ }
\NormalTok{└── course{-}hub/}
\NormalTok{    ├── course{-}hub.Rproj}
\NormalTok{    ├── README.md}
\NormalTok{    ├── week{-}01/}
\NormalTok{    │   ├── tutorials/}
\NormalTok{    │   └── exercises/}
\NormalTok{    └── week{-}02/}
\end{Highlighting}
\end{Shaded}

A trailing \texttt{/} indicates that the item is a \textbf{folder}. In
the above folder structure, skill-building-in-R/, course-hub/, week-01/,
and week-02/ are all folders, not files. Files will have an extension,
like \texttt{README.md}.

\texttt{course-hub.Rproj} is an Rproject file, something we will learn
about next week.

\texttt{README.md} is a file - every repository (\texttt{course-hub})
should have a \texttt{README} file, which states what the repository
holds. The extension \texttt{.md} just means it is a markdown file. A
markdown file is a plain text file that uses a human-readable syntax to
format text. We will learn more about markdown (and its newer cousin,
quarto) next week.

In the \texttt{week-01/} and \texttt{week-02/} folders (and for every
week), there are subfolders for the tutorial templates and the exercises
templates.

In \texttt{course-hub}, you will:

\begin{itemize}
\tightlist
\item
  clone the repository once;
\item
  pull updates before class;
\item
  view course materials; and
\item
  copy the tutorial and exercise files you need.
\end{itemize}

You will \textbf{not} complete or save your work in \texttt{course-hub},
and you do not have the ability to push changes to it.

\subsection{Personal tutorial
repository}\label{personal-tutorial-repository}

You will create one personal repository for all in-class tutorials and
exercises during the semester. This will become a foundation of code you
can refer back to later.

You can name your tutorials-exercise repo:

\begin{Shaded}
\begin{Highlighting}[]
\NormalTok{tutorials{-}exercises}
\end{Highlighting}
\end{Shaded}

Your repository might eventually look like this:

\begin{Shaded}
\begin{Highlighting}[]
\NormalTok{skill{-}building{-}in{-}R/}
\NormalTok{└── tutorials{-}exercises/}
\NormalTok{    ├── README.md}
\NormalTok{    ├── week{-}01/}
\NormalTok{    │   ├── tutorials/}
\NormalTok{    │   └── exercises/}
\NormalTok{    ├── week{-}02/}
\NormalTok{    │   ├── tutorials/}
\NormalTok{    │   └── exercises/}
\NormalTok{    └── tutorials{-}exercises.Rproj}
\end{Highlighting}
\end{Shaded}

In your personal tutorial repository, you will:

\begin{itemize}
\tightlist
\item
  copy tutorial and exercise templates from \texttt{course-hub};
\item
  complete tutorials and exercises;
\item
  save your work;
\item
  commit your changes; and
\item
  push your work to GitHub.
\end{itemize}

I will not be looking at or grading tutorials or exercises. These are
for your own edification. However, your participation points will
reflect your engagement during both.

\begin{center}\rule{0.5\linewidth}{0.5pt}\end{center}

\section{3. Remote and local
repositories}\label{remote-and-local-repositories}

A repository on GitHub is the \textbf{remote} copy. The copy on your
computer is the \textbf{local} copy.

\begin{Shaded}
\begin{Highlighting}[]
\NormalTok{GitHub}
\NormalTok{remote repository stored online}
\NormalTok{        ↓ clone or pull}
\NormalTok{Your computer}
\NormalTok{local repository opened in RStudio}
\NormalTok{        ↑ push}
\end{Highlighting}
\end{Shaded}

Important terms:

{\def\LTcaptype{none} % do not increment counter
\begin{longtable}[]{@{}
  >{\raggedright\arraybackslash}p{(\linewidth - 2\tabcolsep) * \real{0.5000}}
  >{\raggedright\arraybackslash}p{(\linewidth - 2\tabcolsep) * \real{0.5000}}@{}}
\toprule\noalign{}
\begin{minipage}[b]{\linewidth}\raggedright
Term
\end{minipage} & \begin{minipage}[b]{\linewidth}\raggedright
Meaning
\end{minipage} \\
\midrule\noalign{}
\endhead
\bottomrule\noalign{}
\endlastfoot
\textbf{Clone} & Download a repository to your computer for the first
time \\
\textbf{Pull} & Download new changes into an existing local
repository \\
\textbf{Commit} & Save a checkpoint in the repository history on your
computer \\
\textbf{Push} & Upload your local commits to GitHub \\
\end{longtable}
}

You clone a repository only once per computer. After that, you reopen
the existing project and pull updates.

\begin{center}\rule{0.5\linewidth}{0.5pt}\end{center}

\section{4. Organize your course
folders}\label{organize-your-course-folders}

Your repositories should be in the parent folder (next to one another),
not inside one another.

Recommended structure:

\begin{Shaded}
\begin{Highlighting}[]
\NormalTok{skill{-}building{-}in{-}R/}
\NormalTok{├── course{-}hub/}
\NormalTok{├── tutorials{-}exercises/}
\NormalTok{├── homeworks}
\NormalTok{    ├── homework{-}01{-}lastname/}
\NormalTok{    └── homework{-}02{-}lastname/}
\end{Highlighting}
\end{Shaded}

Do \textbf{not} place one Git repository inside another.

Incorrect:

\begin{Shaded}
\begin{Highlighting}[]
\NormalTok{course{-}hub/}
\NormalTok{└── tutorials{-}exercises/}
\end{Highlighting}
\end{Shaded}

Correct:

\begin{Shaded}
\begin{Highlighting}[]
\NormalTok{skill{-}building{-}in{-}R/}
\NormalTok{├── course{-}hub/}
\NormalTok{└── tutorials{-}exercises/}
\end{Highlighting}
\end{Shaded}

\begin{center}\rule{0.5\linewidth}{0.5pt}\end{center}

\section{\texorpdfstring{5. Clone
\texttt{course-hub}}{5. Clone course-hub}}\label{clone-course-hub}

You only need to do this once on each computer you use.

\subsection{On GitHub}\label{on-github}

\begin{enumerate}
\def\labelenumi{\arabic{enumi}.}
\tightlist
\item
  Open the
  \href{https://github.com/fall2026-skill-building-in-R/course-hub}{course-hub
  repository link}
\item
  Click the green \textbf{Code} button.
\item
  Make sure \textbf{HTTPS} is selected.
\item
  Copy the repository URL.
\end{enumerate}

\subsection{In RStudio}\label{in-rstudio}

\begin{enumerate}
\def\labelenumi{\arabic{enumi}.}
\tightlist
\item
  Open RStudio.
\item
  Select \textbf{File → New Project}.
\item
  Select \textbf{Version Control}.
\item
  Select \textbf{Git}.
\item
  Paste the repository URL into \textbf{Repository URL}.
\item
  Choose your \texttt{skill-building-in-R/} parent folder as the project
  directory.
\item
  Click \textbf{Create Project}.
\end{enumerate}

RStudio will clone the repository and open the project. A course-hub
folder will automatically be created through this process.

\subsection{Check your work}\label{check-your-work}

Confirm that:

\begin{itemize}
\tightlist
\item
  \texttt{course-hub} appears as the project name in the upper left-hand
  corner of RStudio;
\item
  the Files pane shows the contents of \texttt{course-hub};
\item
  the Git pane is visible;
\item
  course-hub folder was created in \texttt{skill-building-in-R/} folder;
  and
\item
  you can locate the \texttt{week-01} folder.
\end{itemize}

\begin{center}\rule{0.5\linewidth}{0.5pt}\end{center}

\section{6. Practice pulling an instructor
update}\label{practice-pulling-an-instructor-update}

Your instructor will make a small change to \texttt{course-hub} during
class.

When instructed:

\begin{enumerate}
\def\labelenumi{\arabic{enumi}.}
\tightlist
\item
  Go to
  \href{https://github.com/fall2026-skill-building-in-R/course-hub}{course-hub
  repository link} and note there is no file named \texttt{test.R}
\item
  Make sure the \texttt{course-hub} project is open in RStudio.
\item
  Open the \textbf{Git} pane.
\item
  Click \textbf{Pull}.
\item
  Wait for the pull to finish.
\item
  Locate \texttt{test.R}.
\item
  Go back to
  \href{https://github.com/fall2026-skill-building-in-R/course-hub}{course-hub
  repository link} and note there is now a \texttt{test.R} file.
\end{enumerate}

\subsection{Key idea}\label{key-idea}

\begin{quote}
Clone once. Pull whenever the instructor updates the repository (pull at
the start of class).
\end{quote}

Pulling does not create a second copy of the repository. It updates the
local copy you already have.

\begin{center}\rule{0.5\linewidth}{0.5pt}\end{center}

\section{7. Create your personal tutorial
repository}\label{create-your-personal-tutorial-repository}

Create this repository under your \textbf{personal GitHub account}, not
inside \texttt{course-hub}. This is where you will copy over the
templates for the tutorials and exercises, and complete, save, commit,
and push them.

\subsection{On GitHub}\label{on-github-1}

\begin{enumerate}
\def\labelenumi{\arabic{enumi}.}
\item
  Go to GitHub.
\item
  Click the \textbf{+} menu in the upper-right corner.
\item
  Select \textbf{New repository}.
\item
  Under \textbf{Owner}, select your personal GitHub account (this is
  likely already selected and you will only have options here if you are
  part of organizations).
\item
  Name the repository:

\begin{Shaded}
\begin{Highlighting}[]
\NormalTok{tutorials{-}exercises}
\end{Highlighting}
\end{Shaded}
\item
  Add a short description, such as:

\begin{Shaded}
\begin{Highlighting}[]
\NormalTok{In{-}class tutorials and exercises for Skill Building in R}
\end{Highlighting}
\end{Shaded}
\item
  Choose a `repository visibility' option. The repository visibility is
  asking whether you want this to be a private repo (only you and
  collaborators you add can see it) or public (anyone can see it). I
  usually keep in-progress analyses private and make them public when I
  submit the paper to the journal (so reviewers can see code) or a
  preprint server (so readers can see code). For this class, it does not
  matter as much - choose what you feel like.
\item
  Check \textbf{Add a README file} (this is a description of the
  project, let's look at some other README files).
\item
  Click \textbf{Create repository}.
\end{enumerate}

\begin{center}\rule{0.5\linewidth}{0.5pt}\end{center}

\section{8. Clone your tutorial
repository}\label{clone-your-tutorial-repository}

\subsection{On GitHub}\label{on-github-2}

\begin{enumerate}
\def\labelenumi{\arabic{enumi}.}
\tightlist
\item
  Open your new tutorial repository.
\item
  Click \textbf{Code}.
\item
  Select \textbf{HTTPS}.
\item
  Copy the repository URL.
\end{enumerate}

\subsection{In RStudio}\label{in-rstudio-1}

\begin{enumerate}
\def\labelenumi{\arabic{enumi}.}
\tightlist
\item
  Select \textbf{File → New Project}.
\item
  Select \textbf{Version Control}.
\item
  Select \textbf{Git}.
\item
  Paste the repository URL.
\item
  Choose the same \texttt{skill-building-in-R/} parent folder used for
  \texttt{course-hub}.
\item
  Click \textbf{Create Project}.
\end{enumerate}

Your repositories should now be side by side:

\begin{Shaded}
\begin{Highlighting}[]
\NormalTok{skill{-}building{-}in{-}R/}
\NormalTok{├── course{-}hub/}
\NormalTok{└── tutorials{-}exercises/}
\end{Highlighting}
\end{Shaded}

\begin{center}\rule{0.5\linewidth}{0.5pt}\end{center}

\section{9. Copy the Week 1 materials}\label{copy-the-week-1-materials}

You will complete the tutorial in your personal tutorial repository, not
in \texttt{course-hub} so you will need to copy the tutorial(s) and
exercises for the week into your tutorials-exercises repository.

\subsection{Close or save any open
files}\label{close-or-save-any-open-files}

Before copying files, save your work and close any files you are not
using.

\subsection{Using Finder or File
Explorer}\label{using-finder-or-file-explorer}

\begin{enumerate}
\def\labelenumi{\arabic{enumi}.}
\item
  Open the \texttt{course-hub} folder on your computer.
\item
  Open \texttt{week-01}.
\item
  Copy the \texttt{tutorials} and \texttt{exercises} folders.
\item
  Open your personal \texttt{tutorials-exercises} folder.
\item
  Create a folder named:

\begin{Shaded}
\begin{Highlighting}[]
\NormalTok{week{-}01}
\end{Highlighting}
\end{Shaded}
\item
  Paste the copied materials into \texttt{week-01}.
\end{enumerate}

Your personal repository should now look similar to course-hub, which
looks like this for week-01:

\begin{Shaded}
\begin{Highlighting}[]
\NormalTok{skill{-}building{-}in{-}R/}
\NormalTok{└── tutorials{-}exercises/}
\NormalTok{    ├── README.md}
\NormalTok{    └── week{-}01/}
\NormalTok{        ├── tutorials/}
\NormalTok{        │   └── tutorial{-}01{-}github.md}
\NormalTok{        └── exercises/}
\NormalTok{            └── exercises{-}01{-}github.md}
\end{Highlighting}
\end{Shaded}

\section{10. Open Rstudio as if you are just starting
class}\label{open-rstudio-as-if-you-are-just-starting-class}

When you open RStudio for a working session, you want to open it using
the Rproject file.

To practice this:

\begin{enumerate}
\def\labelenumi{\arabic{enumi}.}
\tightlist
\item
  Quit RStudio so it is not open.
\item
  Navigate to and open your tutorials-exercises/ folder.
\item
  To open this project (and thus RStudio), click on the
  tutorials-exercises.Rproj file.
\item
  Confirm that the project name in the upper left-hand corner of RStudio
  is tutorials-exercises and your files are there in the file pane.
\end{enumerate}

This is how you open RStudio. I always organize my work into RProjects
and so when you click on the .Rproj file, it opens RStudio and your
working directory is set to the folder (\texttt{tutorials-exercises/}).
We will learn more about the \texttt{\{here\}} package next week which
is used in conjunction with Rprojects to solve all working directory
issues.

\section{11. Your weekly in-class
workflow}\label{your-weekly-in-class-workflow}

At the beginning of each class:

\begin{enumerate}
\def\labelenumi{\arabic{enumi}.}
\tightlist
\item
  Open \texttt{course-hub}.
\item
  Pull the latest instructor updates.
\item
  Copy that week's tutorial and exercise materials.
\item
  Paste them into the matching week folder in your
  \texttt{tutorials-exercises} repository.
\end{enumerate}

Then:

\begin{enumerate}
\def\labelenumi{\arabic{enumi}.}
\setcounter{enumi}{4}
\tightlist
\item
  Open your \texttt{tutorials-exercises} RStudio Project.
\item
  Complete the tutorial and exercises.
\item
  Save your files.
\item
  Review your changes in the Git pane.
\item
  Stage the files you want to commit.
\item
  Commit with an informative message.
\item
  Push to GitHub.
\item
  Verify that your work appears in the remote GitHub repository.
\end{enumerate}

\section{12. Repository rules for this
course}\label{repository-rules-for-this-course}

\subsection{\texorpdfstring{In
\texttt{course-hub}}{In course-hub}}\label{in-course-hub}

You should:

\begin{itemize}
\tightlist
\item
  pull updates;
\item
  view course materials; and
\item
  copy starter files.
\end{itemize}

You should \textbf{not}:

\begin{itemize}
\tightlist
\item
  complete your work there;
\item
  commit changes there; or
\item
  push changes there.
\end{itemize}

\subsection{\texorpdfstring{In your \texttt{tutorials-exercises}
repository}{In your tutorials-exercises repository}}\label{in-your-tutorials-exercises-repository}

You should:

\begin{itemize}
\tightlist
\item
  complete tutorials and exercises;
\item
  save your work;
\item
  review your changes;
\item
  commit meaningful checkpoints; and
\item
  push regularly.
\end{itemize}

\section{13. One-sentence workflow}\label{one-sentence-workflow}

\begin{quote}
Pull course materials from \texttt{course-hub}, copy them into your
personal \texttt{tutorials-exercises} repository, complete your work
there, and then review, stage, commit, and push your changes to GitHub.
\end{quote}




\end{document}
